% ═══════════════════════════════════════════════════════════════
% MUSTERLÖSUNG: Repaso Examen - Unidad 5 + 6
% Spanisch Klasse 8 | 40 BE
% ═══════════════════════════════════════════════════════════════

\documentclass[11pt, a4paper]{article}

\usepackage[utf8]{inputenc}
\usepackage[T1]{fontenc}
\usepackage[ngerman]{babel}
\usepackage{helvet}
\renewcommand{\familydefault}{\sfdefault}

\usepackage[margin=1.8cm, top=2cm, bottom=1.5cm]{geometry}
\usepackage{enumitem}
\usepackage{tabularx}
\usepackage{booktabs}
\usepackage{array}
\usepackage{multicol}
\usepackage{amssymb}
\usepackage{eurosym}

\usepackage{xcolor}
\usepackage{tcolorbox}
\tcbuselibrary{skins}

\usepackage{fancyhdr}
\usepackage{lastpage}

% ═══════════════════════════════════════════════════════════════
% FARBEN
% ═══════════════════════════════════════════════════════════════

\definecolor{spanisch}{HTML}{c41e3a}
\definecolor{grau}{HTML}{888888}
\definecolor{hellgrau}{HTML}{f5f5f5}
\definecolor{loesung}{HTML}{c41e3a}

% ═══════════════════════════════════════════════════════════════
% VARIABLEN
% ═══════════════════════════════════════════════════════════════

\newcommand{\thema}{Examen: Unidad 5 + 6}
\newcommand{\fach}{Spanisch}
\newcommand{\klasse}{8}
\newcommand{\maxpunkte}{40}
\newcommand{\zeit}{45 Minuten}

% ═══════════════════════════════════════════════════════════════
% KOPF- UND FUSSZEILE
% ═══════════════════════════════════════════════════════════════

\pagestyle{fancy}
\fancyhf{}
\renewcommand{\headrulewidth}{0.4pt}
\renewcommand{\footrulewidth}{0pt}

\lhead{\small\fach{} Kl. \klasse}
\chead{\small\textcolor{loesung}{\textbf{MUSTERLÖSUNG}}}
\rhead{\small Seite \thepage/\pageref{LastPage}}

% ═══════════════════════════════════════════════════════════════
% BEFEHLE
% ═══════════════════════════════════════════════════════════════

\newcommand{\punktebox}[1]{\hfill\fbox{\scriptsize \_\_/#1}}
\newcommand{\luecke}[1]{\rule{#1}{0.4pt}}
\newcommand{\es}[1]{\textcolor{spanisch}{\textbf{#1}}}
\newcommand{\lsg}[1]{\textcolor{loesung}{\textbf{#1}}}

% Bilinguale Aufgabenstellung
\newcommand{\aufgabe}[3]{%
    \textbf{\textcolor{spanisch}{Ejercicio #1: #2}} \punktebox{#3}\\[2pt]
}
\newcommand{\enunciado}[2]{%
    \textbf{#1}\\[1pt]
    {\small\textit{#2}}%
}

\newtcolorbox{textbox}{
    colback=hellgrau,
    colframe=grau,
    boxrule=0.5pt,
    arc=0pt,
    left=8pt, right=8pt, top=6pt, bottom=6pt
}

\setlength{\parindent}{0pt}
\setlength{\parskip}{6pt}

% ═══════════════════════════════════════════════════════════════
% DOKUMENT
% ═══════════════════════════════════════════════════════════════

\begin{document}

% ═══════════════════════════════════════════════════════════════
% KOPFBEREICH
% ═══════════════════════════════════════════════════════════════

\begin{center}
    {\Large\bfseries\textcolor{spanisch}{Examen: Repaso Unidad 5 + 6}}\\[4pt]
    {\small Mi ropa preferida -- Mi lugar preferido}\\[4pt]
    {\large\textcolor{loesung}{\textbf{--- MUSTERLÖSUNG ---}}}
\end{center}

\vspace{4pt}

\begin{tabularx}{\textwidth}{|X|l|l|}
\hline
\textbf{Nombre:} \lsg{Musterlösung} &
\textbf{Clase:} \klasse &
\textbf{Fecha:} \lsg{---} \\
\hline
\end{tabularx}

\vspace{4pt}

\begin{tabularx}{\textwidth}{|X|X|}
\hline
\textbf{Tiempo:} \zeit &
\textbf{Puntuación:} \lsg{40} / \maxpunkte{} BE \\
\hline
\end{tabularx}

\vspace{6pt}

% ═══════════════════════════════════════════════════════════════
% BEWERTUNGSMATRIX (14-NP Sek I) - 40 BE
% ═══════════════════════════════════════════════════════════════

\begin{center}
\footnotesize
\begin{tabular}{|c|c|c|c|c|c|c|c|c|c|c|c|c|c|c|c|}
\hline
\textbf{NP} & 14 & 13 & 12 & 11 & 10 & 9 & 8 & 7 & 6 & 5 & 4 & 3 & 2 & 1 & 0 \\
\hline
\textbf{BE} & 40 & 38 & 36 & 34 & 32 & 29 & 27 & 24 & 21 & 19 & 16 & 13 & 11 & 8 & 0 \\
\hline
\end{tabular}
\end{center}

\vspace{8pt}

% ═══════════════════════════════════════════════════════════════
% EJERCICIO 1: HÖRVERSTEHEN (4 BE)
% ═══════════════════════════════════════════════════════════════

\aufgabe{1}{Comprensión auditiva -- Los precios}{4}

\enunciado{Escucha y escribe los precios en números.}{Höre zu und schreibe die Preise als Ziffern.}

\vspace{4pt}
{\small\textcolor{grau}{[Audio: Pista 17 -- Beispielpreise]}}

\vspace{6pt}

\begin{tabular}{@{}p{6cm}p{4cm}p{2cm}@{}}
a) La camiseta cuesta & \lsg{125} \euro & (1 BE) \\[8pt]
b) El jersey cuesta & \lsg{249} \euro & (1 BE) \\[8pt]
c) Los pantalones cuestan & \lsg{350} \euro & (1 BE) \\[8pt]
d) La chaqueta cuesta & \lsg{575} \euro & (1 BE) \\
\end{tabular}

\vspace{4pt}
{\footnotesize\textcolor{grau}{Hinweis: Die konkreten Zahlen hängen vom Audio ab.}}

\vspace{10pt}

% ═══════════════════════════════════════════════════════════════
% EJERCICIO 2: VOCABULARIO - LA ROPA (4 BE)
% ═══════════════════════════════════════════════════════════════

\aufgabe{2}{Vocabulario -- La ropa}{4}

\enunciado{Escribe el nombre de cada prenda con el artículo.}{Schreibe den Namen jedes Kleidungsstücks mit Artikel.}

\vspace{6pt}

\begin{tabular}{@{}p{2.5cm}p{4cm}p{2.5cm}p{4cm}p{2cm}@{}}
a) Pullover: & \lsg{el jersey} & b) T-Shirt: & \lsg{la camiseta} & (je 1 BE) \\[8pt]
c) Hose: & \lsg{los pantalones} & d) Rock: & \lsg{la falda} & \\
\end{tabular}

\vspace{10pt}

% ═══════════════════════════════════════════════════════════════
% EJERCICIO 3: LOS COLORES + ADJEKTIVE (4 BE)
% ═══════════════════════════════════════════════════════════════

\aufgabe{3}{Los colores -- Adjektivanpassung}{4}

\enunciado{Completa con el color correcto. ¡Atención a la terminación!}{Ergänze mit der richtigen Farbe. Achte auf die Endung!}

\vspace{6pt}

\begin{tabular}{@{}p{9cm}p{3cm}p{2cm}@{}}
a) La camiseta es \lsg{blanca}. & (weiß) & (1 BE) \\[8pt]
b) El jersey es \lsg{negro}. & (schwarz) & (1 BE) \\[8pt]
c) Los zapatos son \lsg{marrones}. & (braun) & (1 BE) \\[8pt]
d) La falda es \lsg{roja}. & (rot) & (1 BE) \\
\end{tabular}

\vspace{10pt}

\newpage

% ═══════════════════════════════════════════════════════════════
% EJERCICIO 4: VERBOS MODALES (4 BE)
% ═══════════════════════════════════════════════════════════════

\aufgabe{4}{Verbos modales -- querer, poder, tener que}{4}

\enunciado{Completa con la forma correcta del verbo entre paréntesis.}{Ergänze mit der richtigen Form des Verbs in Klammern.}

\vspace{6pt}

\begin{tabular}{@{}p{10cm}p{3cm}p{2cm}@{}}
a) Yo \lsg{quiero} comprar un jersey nuevo. & (querer) & (1 BE) \\[8pt]
b) ¿Tú \lsg{puedes} ayudarme? & (poder) & (1 BE) \\[8pt]
c) Nosotros \lsg{tenemos que} estudiar para el examen. & (tener que) & (1 BE) \\[8pt]
d) Vosotros \lsg{tenéis que} hacer los deberes. & (tener que) & (1 BE) \\
\end{tabular}

\vspace{10pt}

% ═══════════════════════════════════════════════════════════════
% EJERCICIO 5: VERBO IR + A/AL (4 BE)
% ═══════════════════════════════════════════════════════════════

\aufgabe{5}{Verbo ir + a / al}{4}

\enunciado{Completa con la forma correcta de \emph{ir} y \emph{a} o \emph{al}.}{Ergänze mit der richtigen Form von \emph{ir} und \emph{a} oder \emph{al}.}

\vspace{6pt}

\begin{tabular}{@{}p{12cm}p{2cm}@{}}
a) Yo \lsg{voy} \lsg{al} cine con mis amigos. & (1 BE) \\[8pt]
b) María \lsg{va} \lsg{a la} tienda de ropa. & (1 BE) \\[8pt]
c) Nosotros \lsg{vamos} \lsg{al} centro comercial. & (1 BE) \\[8pt]
d) ¿Vosotros \lsg{vais} \lsg{a la} playa mañana? & (1 BE) \\
\end{tabular}

\vspace{10pt}

% ═══════════════════════════════════════════════════════════════
% EJERCICIO 6: LOS NÚMEROS (4 BE)
% ═══════════════════════════════════════════════════════════════

\aufgabe{6}{Los números 100-1000}{4}

\enunciado{Escribe los números como palabras españolas.}{Schreibe die Zahlen als spanische Wörter.}

\vspace{6pt}

\begin{tabular}{@{}p{2cm}p{5.5cm}p{2cm}p{5.5cm}p{2cm}@{}}
a) 250 = & \lsg{doscientos cincuenta} & b) 500 = & \lsg{quinientos} & (je 1 BE) \\[8pt]
c) 715 = & \lsg{setecientos quince} & d) 999 = & \lsg{novecientos noventa y nueve} & \\
\end{tabular}

\vspace{10pt}

% ═══════════════════════════════════════════════════════════════
% EJERCICIO 7: EL TIEMPO (4 BE)
% ═══════════════════════════════════════════════════════════════

\aufgabe{7}{El tiempo -- Das Wetter}{4}

\enunciado{¿Qué tiempo hace? Escribe la expresión correcta.}{Wie ist das Wetter? Schreibe den richtigen Ausdruck.}

\vspace{6pt}

\begin{tabular}{@{}p{5cm}p{6cm}p{2cm}@{}}
a) Es ist sonnig: & \lsg{Hace sol.} & (1 BE) \\[8pt]
b) Es regnet: & \lsg{Llueve.} & (1 BE) \\[8pt]
c) Es ist kalt: & \lsg{Hace frío.} & (1 BE) \\[8pt]
d) Es ist bewölkt: & \lsg{Está nublado.} & (1 BE) \\
\end{tabular}

\vspace{10pt}

\newpage

% ═══════════════════════════════════════════════════════════════
% EJERCICIO 8: LESEVERSTEHEN (4 BE)
% ═══════════════════════════════════════════════════════════════

\aufgabe{8}{Comprensión lectora}{4}

\enunciado{Lee el texto y contesta a las preguntas en español.}{Lies den Text und beantworte die Fragen auf Spanisch.}

\vspace{6pt}

\begin{textbox}
\small
Hoy es sábado y Lucía va al centro comercial con su madre. Lucía quiere comprar ropa nueva porque tiene una fiesta el próximo fin de semana. En la tienda, Lucía ve un vestido rojo muy bonito. El vestido cuesta doscientos cincuenta euros, pero es muy caro. Después, Lucía ve una falda azul que cuesta ochenta y cinco euros. La falda le queda muy bien. Al final, Lucía compra la falda azul y una camiseta blanca. Está muy contenta.
\end{textbox}

\vspace{6pt}

\begin{tabular}{@{}p{0.5cm}p{12cm}p{2cm}@{}}
a) & ¿Adónde va Lucía el sábado? & (1 BE) \\[4pt]
& \lsg{Lucía va al centro comercial (con su madre).} & \\[8pt]
b) & ¿Por qué quiere comprar ropa nueva? & (1 BE) \\[4pt]
& \lsg{Porque tiene una fiesta el próximo fin de semana.} & \\[8pt]
c) & ¿Cuánto cuesta el vestido rojo? & (1 BE) \\[4pt]
& \lsg{El vestido cuesta doscientos cincuenta euros. / 250 euros.} & \\[8pt]
d) & ¿Qué compra Lucía al final? & (1 BE) \\[4pt]
& \lsg{Compra una falda azul y una camiseta blanca.} & \\
\end{tabular}

\vspace{10pt}

% ═══════════════════════════════════════════════════════════════
% EJERCICIO 9: EXPRESIÓN ESCRITA (8 BE)
% ═══════════════════════════════════════════════════════════════

\aufgabe{9}{Expresión escrita -- Diálogo en una tienda}{8}

\enunciado{Escribe un diálogo en una tienda de ropa. Usa las expresiones del recuadro.}{Schreibe einen Dialog in einem Kleidungsgeschäft. Benutze die Ausdrücke aus dem Kasten.}

\vspace{4pt}

\textbf{Situación:} Quieres comprar una camiseta verde. \textit{(Du möchtest ein grünes T-Shirt kaufen.)}

\vspace{6pt}

\textcolor{loesung}{\textbf{Beispiellösung:}}

\vspace{4pt}

\begin{tabular}{@{}p{2.2cm}p{12.5cm}@{}}
\textbf{Vendedor:} & \lsg{¡Hola! ¿En qué puedo ayudarte?} \\[8pt]
\textbf{Cliente:} & \lsg{Hola. Quiero comprar una camiseta verde.} \\[8pt]
\textbf{Vendedor:} & \lsg{Muy bien. ¿Qué talla tienes?} \\[8pt]
\textbf{Cliente:} & \lsg{Tengo la talla M. ¿Cuánto cuesta?} \\[8pt]
\textbf{Vendedor:} & \lsg{Cuesta treinta y cinco euros.} \\[8pt]
\textbf{Cliente:} & \lsg{Perfecto, me queda bien. Me la llevo.} \\
\end{tabular}

\vspace{8pt}

{\small\textcolor{grau}{\textbf{Bewertungskriterien (je 2 BE):}
\begin{itemize}[leftmargin=15pt, itemsep=1pt]
    \item Verwendung der Dialogausdrücke (2 BE)
    \item Korrekte Verbformen (2 BE)
    \item Logischer Dialogaufbau (2 BE)
    \item Sprachliche Korrektheit (2 BE)
\end{itemize}}}

\vspace{10pt}

{\footnotesize\textcolor{grau}{\textbf{Audio Ejercicio 1:} Pista 17 (basiert auf Übung 16a, Qué pasa 1, S. 98)}}

\end{document}
